% Zone „Kennst du schon" – Hauptnummern 1 bis 11
\einheitenkopf[zone][Kennst du schon]{Das kennst du schon}
\setcounter{aufgabe}{0}

\begin{aufgabe}{Ich kann Zahlen quadrieren, auch negative Zahlen und Dezimalzahlen.}
Berechne.
\begin{teilezwei}
\tz $6^2 = {}$ \leerfeld & \tz $10^2 = {}$ \leerfeld \\
\tz $(-8)^2 = {}$ \leerfeld & \tz $2{,}5^2 = {}$ \leerfeld \\
\tz $(-0{,}5)^2 = {}$ \leerfeld & \\
\end{teilezwei}
\end{aufgabe}

\begin{aufgabe}{Ich kann mit Klammern und negativen Zahlen rechnen: Klammer zuerst, Punkt vor Strich.}
Berechne.
\begin{teilezwei}
\tz $3 \cdot (2 + 4) = {}$ \leerfeld & \tz $20 - 4 \cdot 3 = {}$ \leerfeld \\
\tz $(-2) \cdot 6 = {}$ \leerfeld & \tz $(-4) \cdot (-4 + 9) = {}$ \leerfeld \\
\tz $2 \cdot (-3)^2 - 10 = {}$ \leerfeld & \\
\end{teilezwei}
\end{aufgabe}

\begin{aufgabe}{Ich kann eine lineare Gleichung lösen.}
Löse die Gleichung.
\begin{gleichungsraster}[2]
\gl{x + 4 = 10} & \gl{3x = 12} \\
\gl{x - 2{,}5 = 4} & \gl{x + 6 = 0} \\
\gl{2x - 5 = 11} & \gl{-x = 7} \\
\end{gleichungsraster}
\end{aufgabe}

\begin{aufgabe}{Ich kann durch Einsetzen prüfen, ob eine Zahl eine Lösung ist.}
Setze die Zahl ein und kreuze an: wahre Aussage (wA) oder falsche Aussage (fA).
\begin{teile}
\teil $x = 3$ \quad in \quad $x + 5 = 8$ \hfill \kreuz{wA} \kreuz{fA}
\teil $x = 2$ \quad in \quad $4x = 10$ \hfill \kreuz{wA} \kreuz{fA}
\teil $x = 0{,}5$ \quad in \quad $4x + 1 = 3$ \hfill \kreuz{wA} \kreuz{fA}
\teil $x = -2$ \quad in \quad $5 - x = 3$ \hfill \kreuz{wA} \kreuz{fA}
\teil $x = -3$ \quad in \quad $2x - 1 = -7$ \hfill \kreuz{wA} \kreuz{fA}
\end{teile}
\end{aufgabe}

\begin{aufgabe}{Ich kann Quadratwurzeln ziehen.}
Gib die Wurzel an. Gibt es sie nicht, schreibe „gibt es nicht“.
\begin{teilezwei}
\tz $\sqrt{49} = {}$ \leerfeld & \tz $\sqrt{100} = {}$ \leerfeld \\
\tz $\sqrt{0{,}64} = {}$ \leerfeld & \tz $\sqrt{196} = {}$ \leerfeld \\
\tz $\sqrt{-9} = {}$ \leerfeld & \\
\anweisung{Rechne mit dem Taschenrechner und runde auf zwei Stellen nach dem Komma.}
\tz $\sqrt{7} \approx {}$ \leerfeld & \\
\end{teilezwei}
\end{aufgabe}

\begin{aufgabe}{Ich kann den Scheitel aus der Scheitelpunktform ablesen.}
Lies den Scheitel S der Parabel ab.
\begin{teilezwei}
\tz $y = (x - 2)^2 + 1$ \quad \punktfeld & \tz $y = (x - 4)^2$ \quad \punktfeld \\
\tz $y = x^2 - 5$ \quad \punktfeld & \tz $y = (x + 3)^2 - 2$ \quad \punktfeld \\
\end{teilezwei}
\end{aufgabe}

\begin{aufgabe}{Ich kann Klammern ausmultiplizieren, auch mit den binomischen Formeln.}
Multipliziere aus.
\begin{gleichungsraster}[1]
\gl{2 \cdot (x + 5)} & \gl{x \cdot (x + 3)} \\
\gl{x \cdot (x - 6)} & \gl{(x + 6)^2} \\
\gl{(x - 9)^2} & \\
\end{gleichungsraster}
\end{aufgabe}

\begin{aufgabe}{Ich kann ausklammern.}
Klammere so viel wie möglich aus.
\begin{gleichungsraster}[1]
\gl{4x + 8} & \gl{x^2 + 3x} \\
\gl{x^2 - 9x} & \gl{2x^2 + 6x} \\
\gl{x^2 - x} & \\
\end{gleichungsraster}
\end{aufgabe}

\begin{aufgabe}{Ich kann Terme ordnen und zusammenfassen.}
Ordne nach $x^2$, $x$ und Zahl und fasse zusammen.
\begin{gleichungsraster}[1]
\gl{3 + x^2 + 2x} & \gl{x^2 + 4x + 3x + 1} \\
\gl{5x - 2 + x^2 - 6} & \gl{2x - x^2 + 4} \\
\gl{x^2 - 3x + 10 - 12 + x} & \\
\end{gleichungsraster}
\end{aufgabe}

\begin{aufgabe}{Ich finde den Fehler beim Auflösen einer quadrierten Klammer.}
Jonas rechnet so:
\rechnung{(x + 5)^2 - 30 &= x^2 + 25 - 30 \\ &= x^2 - 5}
Finde den Fehler, benenne ihn und korrigiere die Rechnung.
\begin{gleichungsraster}[3]
\gl{(x + 5)^2 - 30} & \\
\end{gleichungsraster}
\end{aufgabe}

\begin{aufgabe}{Ich kann eine quadrierte Klammer auflösen und zusammenfassen.}
Multipliziere aus und fasse zusammen.
\begin{gleichungsraster}[2]
\gl{(x + 2)^2 - 7} & \gl{(x - 3)^2 + 4x} \\
\gl{(x - 1)^2 - 2x + 5} & \\
\end{gleichungsraster}
\end{aufgabe}
